% =============================================================================
% From Eclipse Demon to Lunar Apogee
% Companion paper — pdflatex + CJKutf8, two-column
% =============================================================================
\documentclass[10pt,a4paper]{article}

% --- Packages (matches when-the-stars-disagree.tex) -------------------------
\usepackage[utf8]{inputenc}
\usepackage[T1]{fontenc}
\usepackage{CJKutf8}
\usepackage{amsmath,amssymb,amsfonts}
\usepackage{booktabs}
\usepackage{multirow}
\usepackage{graphicx}
\usepackage{xcolor}
\usepackage{hyperref}
\usepackage[margin=0.75in]{geometry}
\usepackage{enumitem}
\usepackage{float}
\usepackage{caption}
\usepackage{subcaption}
\usepackage{siunitx}
\usepackage{longtable}
\usepackage[numbers,sort&compress]{natbib}
\usepackage{xurl}
\usepackage{authblk}
\usepackage{pgfplots}
\usepackage{tikz}
\usepackage{gensymb}
\usepackage{adjustbox}
\usepackage[framemethod=tikz]{mdframed}

\pgfplotsset{compat=1.18}

\setlength{\columnsep}{18pt}
\setlength{\emergencystretch}{3em}

\hypersetup{
    colorlinks=true,
    linkcolor=blue!70!black,
    citecolor=green!50!black,
    urlcolor=blue!60!black,
}

% CJK helper
\newcommand{\zh}[1]{\begin{CJK}{UTF8}{bsmi}#1\end{CJK}}

% ORCID helper
\newcommand{\orcidlink}[1]{%
  \href{https://orcid.org/#1}{\textcolor{green!50!black}{\textsf{iD}}}%
}

% pgfplots colors
\definecolor{apogeeblue}{HTML}{2166AC}
\definecolor{nodered}{HTML}{D6604D}

% --- Title -------------------------------------------------------------------

\title{%
  \textbf{From Eclipse Demon to Lunar Apogee:}\\[4pt]
  \large Six Modes of Transformation in Cross-Civilizational\\
  Astronomical Transmission%
}

\author[1]{Albert Kwun Tai Hui\,\orcidlink{0009-0000-8978-0231}}
\affil[1]{Security Ronin}

\date{March 2026}

% =============================================================================
\begin{document}

\twocolumn[
\maketitle

\begin{abstract}
When astronomical concepts travel across civilizational boundaries, they do not
arrive unchanged. This paper argues that such transformations are systematic
rather than random, and proposes a six-mode theoretical framework---transliteration,
reidentification, indigenous restructuring, hybridization, foreign correction, and
algorithmic fossilization---derived from an extended case study spanning thirteen
centuries.

The case study follows the Indian shadow planet Ketu (\zh{計都}) and its coordinate
reference frame as both were absorbed into Chinese \zh{七政四餘} (Seven Governors,
Four Remainders) astrology.  On the coordinate axis the paper traces four
successive regime changes: ecliptic sidereal coordinates in the Song dynasty
(\zh{三辰通載}); equatorial sidereal in the Yuan (\zh{鄭氏星案}); a novel
quasi-ecliptic sidereal hybrid in the Ming (\zh{果老星宗}); and ecliptic tropical
under the Jesuit reform of 1645.  In parallel, Ketu's astronomical identity was
quietly reidentified---from the descending lunar node to the osculating lunar
apogee---a shift demonstrated by Niu Weixing's 1994 recomputation of Tang-dynasty
ephemeris data \cite{niu1995ketu}, before being reverted by Schall von Bell in the
\textit{Shixian} calendar.

The paper contributes: (1)~computational verification of the apsidal identification
using modern JPL ephemeris algorithms; (2)~the first rigorous definition of the
\zh{似黃道} (quasi-ecliptic) coordinate system; and (3)~the six-mode framework as
a reusable instrument for the comparative history of transmitted astronomy.
\end{abstract}
\vspace{12pt}
]

% === §1 Introduction =========================================================
\section{Introduction}\label{sec:intro}
% ~1,500 words — Task 6

\subsection*{The Silk Road as a Channel for Astronomical Ideas}

Astronomical knowledge traveled the Silk Road alongside silk, spices, and
religious texts.  Between the second and ninth centuries~CE, Greek planetary
models, Indian \textit{siddh\=anta} tables, Persian \textit{z\={\i}j} manuals,
and Babylonian omen compendia all made their way into Central Asia and China via
a network of court astronomers, Buddhist missionaries, Sogdian merchants, and
imperial tribute missions.  Unlike a bolt of silk, however, an astronomical concept
does not arrive at its destination intact.  It is translated into a new language,
fitted into an existing institutional structure, combined with native techniques,
challenged by fresh observations, and eventually enshrined in software whose users
have forgotten what it ever meant.  These transformations leave detectable traces
in historical texts---and in the computed positions of celestial bodies.

This paper argues that the transformations are systematic rather than random.
Using a single case study that spans thirteen centuries and four Chinese dynasties,
we identify six recurring patterns---six \emph{modes} of transformation---that
structured the reception of Indian astronomy in China.  The case study concerns
two related objects: the Indian shadow planet Ketu (\zh{計都}) and the coordinate
reference frame used to specify its position.  Both entered Chinese astronomy in
718~CE through the translation of the \textit{Jiuzhi li} (\zh{九執曆},
\textit{Nine Luminaries Calendar}) by Gautama Siddh\=artha (\zh{瞿曇悉達}), a
scholar of Indian descent serving at the Tang court.  Both underwent a remarkable
sequence of transformations over the following millennium, and both continue to
influence popular Chinese-language astrology software today.

\subsection*{Two Parallel Threads}

The paper traces two intertwined threads that run from 718~CE to the present.

\textit{Thread~1: The coordinate reference frame.}  Every astronomical position
requires a reference frame---a choice of origin, fundamental plane, and direction
of increasing longitude.  The original Indian system used ecliptic coordinates
measured from a sidereal zero point (a fixed star rather than the moving
vernal equinox).  Chinese astronomy in the Tang and Song dynasties preferred
equatorial coordinates referenced to the twenty-eight lunar lodges (\zh{二十八宿},
\textit{xiu}).  When Ketu entered the Chinese tradition, a tension arose between
the imported ecliptic frame and the native equatorial one.  That tension was
resolved---and re-resolved---four times:
\begin{enumerate}[leftmargin=*, label=(\arabic*)]
  \item \textbf{Song ecliptic sidereal.}  The \zh{三辰通載} (c.~11th century~CE)
    presents Ketu positions in a coordinate frame that is recognisably ecliptic
    and sidereal, close to the Indian original.
  \item \textbf{Yuan equatorial sidereal.}  The \zh{鄭氏星案} (late Yuan, c.~14th
    century) converts all positions into equatorial coordinates aligned with the
    native lodge system---a thoroughgoing indigenization of the imported framework.
  \item \textbf{Ming quasi-ecliptic sidereal.}  The \zh{果老星宗} (Ming dynasty,
    c.~15th--16th century) introduces what we term the \zh{似黃道} (quasi-ecliptic)
    frame: a hybrid that applies ecliptic-style longitude increments along a plane
    that is neither the true ecliptic nor the equator.  This is the least-studied
    coordinate system in Chinese historical astronomy, and the paper offers the
    first formal definition.
  \item \textbf{Qing ecliptic tropical.}  The Jesuit-led calendar reform of 1645
    (\textit{Shixian li}, \zh{時憲曆}) reimposed European tropical ecliptic
    coordinates, resetting the accumulated drift and discarding the Ming hybrid.
\end{enumerate}
The trajectory is not a smooth evolution but a sequence of reversals---``there
and back again, but not quite'': the Qing frame resembles the Song one in being
ecliptic, but has shifted from sidereal to tropical reckoning.

\textit{Thread~2: The identity of Ketu.}  In Indian astronomy Ketu (\textit{ketu},
literally ``comet'' or ``banner'') designates the descending lunar node---one of
the two points where the Moon's orbit crosses the ecliptic.  When the Moon crosses
the descending node moving southward, conditions favour a lunar eclipse; when
it crosses the ascending node Rahu (\zh{羅睺}), conditions favour a solar eclipse.
The two nodes are therefore deeply embedded in eclipse prediction.

Niu Weixing's landmark 1994 analysis of the \zh{七曜攘災訣} (806~CE), a Tang
ephemeris attributed to the monk Amoghavajra's circle, demonstrated that the
Ketu positions recorded therein do \emph{not} track the descending node
\cite{niu1995ketu}.  They track the osculating lunar apogee---the point in the
Moon's elliptical orbit farthest from Earth.  The apogee was known to Indian
astronomers as \textit{mand\=occa} (``slow apex'') and in Persian-influenced
sources as \textit{K\=etus}; its conflation with the node appears to have
occurred during or shortly after the Tang transmission.  By the Ming period,
Ketu as apogee was thoroughly assimilated into \zh{七政四餘} practice, where it
joined Rahu (\zh{羅睺}), the Purple Star (\zh{紫氣}), and the Monthly Comet
(\zh{月孛}) as the four ``remainders''---invisible bodies that governed fate
alongside the seven visible planets.  The Jesuit reform then reverted Ketu to
the descending node, erasing five centuries of accumulated reidentification.

\subsection*{Six Modes of Transformation}

The six modes proposed here are not a priori categories imposed on the data.
They emerge from the case study as the recurring patterns by which astronomical
concepts were altered during transmission.  We state them here in condensed form;
each receives fuller treatment in the sections that follow.

\textbf{Mode~1: Transliteration.}  Foreign technical terms are adopted into the
receiving tradition with minimal semantic change, using phonetic approximation.
\textit{R\=ahu} becomes \zh{羅睺} (\textit{lu\'ohou}); \textit{Ketu} becomes
\zh{計都} (\textit{j\`{\i}d\=u}).  The sounds are borrowed; the astronomical
meanings---initially---travel with them.  Transliteration is the entry mechanism:
it creates a named slot in the receiving system before that slot has been filled
with indigenous content.

\textbf{Mode~2: Reidentification.}  The name persists while the astronomical
referent shifts.  The transformation is invisible at the lexical level but
detectable by comparing computed positions against historical ephemeris records.
Ketu's drift from descending node to osculating apogee is the central instance;
we suspect the shift happened in stages between 718 and 806~CE as Tang
practitioners noticed that the imported node-based algorithms gave increasingly
poor results and sought a better-fitting celestial body.

\textbf{Mode~3: Indigenous restructuring.}  The native tradition actively
overwrites the imported framework by substituting its own coordinate system,
calendar reckoning, or theoretical structure.  The Yuan equatorial conversion
is the clearest example: Yuan astronomers did not argue with the Indian ecliptic
frame; they simply replaced it with the equatorial lodge-coordinate system that
their \textit{Shoushi li} (\zh{授時曆}) reform had just made canonical.

\textbf{Mode~4: Hybridization.}  Elements of the imported and native systems are
combined into a novel structure that belongs fully to neither parent tradition.
The Ming \zh{似黃道} system is the prime instance.  It preserves ecliptic-style
longitude from the Indian source while anchoring the zero-point and plane to a
native star-count system; the result is not Indian ecliptic astronomy and not
Chinese lodge astronomy, but a tertium quid that persists in Chinese-language
software to this day.

\textbf{Mode~5: Foreign correction.}  An external authority---a new wave of
transmission, a conquering dynasty, or an imported expertise---imposes a reset
that overrides accumulated indigenous development.  Schall von Bell's 1645
\textit{Shixian li} reform is the paradigm case.  Backed by Qing imperial
authority, it discarded the Ming quasi-ecliptic frame, reverted Ketu to the
descending node, and replaced sidereal with tropical reckoning---all in a single
edict.

\textbf{Mode~6: Algorithmic fossilization.}  A living astronomical practice,
once encoded in tables or (later) software, becomes separated from its theoretical
foundation.  Users apply the algorithm without access to---or interest in---the
historical reasoning behind it.  Errors accumulate unnoticed because the algorithm
is self-consistent even if its foundations have shifted.  This mode is not
historically confined to the past: modern Chinese-language astrology software
routinely inherits the Ming \zh{似黃道} frame without labelling it as such,
producing systematic positional errors of several degrees when compared against
modern ephemerides.

\subsection*{Method}

The paper combines close reading of primary astronomical texts with computational
verification using modern ephemeris algorithms.  For the computational component
we use the Swiss Ephemeris library (version~2.10), which implements the JPL
DE441 planetary and lunar solution.  This allows us to ask a precise empirical
question: does a given historical Ketu ephemeris track the lunar apogee, the
descending node, or something else?  We convert historical sexagesimal positions
to modern decimal degrees, apply the appropriate calendar and epoch corrections,
and compare against Swiss Ephemeris output for the descending node and for the
osculating apogee at the same Universal Time.  The method follows and extends
the approach used by Niu Weixing \cite{niu1995ketu} for the \zh{七曜攘災訣} and
applied to Ming texts by Niu in his later synthesis \cite{niu2010fourremains}.
Computational details and residual tables are presented in \S\ref{sec:computation}.

\subsection*{Relationship to Existing Scholarship}

This paper builds on four converging research programmes.

Niu Weixing's identification of Ketu as the lunar apogee
\cite{niu1995ketu,niu2010fourremains} is the single most important prior result
for our argument.  Without Niu's recomputation, the reidentification mode would
remain invisible.  Our contribution is to situate Niu's finding within a broader
transformative sequence and to provide independent computational confirmation
using a modern ephemeris unavailable to Niu in 1994.

Bill Mak's work on pseudo-planets traces the historical trajectory of Rahu and
Ketu from Hellenistic and Indian origins into China and Japan, with particular
attention to Persian-mediated transmission \cite{mak2014pseudoplanets,mak2024persian}.
Mak documents the textual genealogy that brought shadow-planet concepts to East
Asia; the present paper takes up where Mak's textual history ends, asking what
happened to those concepts inside China over the long Ming--Qing transition.

Jeffrey Kotyk has produced detailed studies of the sinicization of Indo-Iranian
astrology in the Tang and Song courts \cite{kotyk2018sinicization,kotyk2017buddhist}.
Kotyk establishes the Buddhist monastic networks---particularly Amoghavajra's
school---through which Indian astronomical texts circulated in the eighth century.
Our \S\ref{sec:tang} on the Tang transmission depends heavily on Kotyk's
reconstruction of institutional context.

Liu Ts'un-yan's work on the Jesuit encounter with Chinese astrology
\cite{liu2020missionary} provides the documentary foundation for our account of
Mode~5 (foreign correction) in \S\ref{sec:jesuit}.  The Jesuit reform was not
merely a technical event; it was an act of epistemic authority, and Liu's analysis
of Schall von Bell's rhetorical strategies illuminates why the reversion of Ketu
to the descending node was politically possible in 1645.

A companion paper, ``When the Stars Disagree: A Cross-Civilizational Analysis
of Computational Errors in Astronomical Divination,'' presents a broader taxonomy
of error types arising from coordinate-frame mismatches across the Chinese,
Hellenistic, Indian, and Persian traditions.  The present paper concentrates on
the transformation \emph{mechanisms} rather than the error \emph{outcomes}.

\subsection*{Structure of the Paper}

Section~\ref{sec:indian} describes the Indian original: the conceptual status of
Rahu and Ketu in Sanskrit \textit{siddh\=anta} astronomy and the ecliptic sidereal
coordinate frame in which they were embedded.  Sections~\ref{sec:tang}--\ref{sec:jesuit}
trace the four dynasties and four coordinate regimes in sequence.
Section~\ref{sec:fossil} analyses the ongoing algorithmic fossilization in modern
software.  Section~\ref{sec:computation} presents the computational verification.
Section~\ref{sec:discussion} draws the six modes together into the theoretical
framework and discusses their applicability beyond the Ketu case.
Section~\ref{sec:conclusion} concludes.

% === §2 The Indian Original ===================================================
\section{The Indian Original: Rahu, Ketu, and Ecliptic Sidereal Coordinates}
\label{sec:indian}
% ~1,000 words — Task 7

\subsection*{Ecliptic Sidereal Coordinates}

Indian \textit{Jyotish} astronomy rests on a coordinate framework that differs
from both the Greek tropical system and the native Chinese equatorial tradition.
The fundamental reference plane is the \emph{ecliptic}---the apparent annual path
of the Sun against the background of fixed stars.  Along this plane the sky is
divided into twelve \textit{r\=asi} (zodiac signs), each occupying exactly 30° of
arc, yielding the familiar circle of 360°.  What distinguishes the Indian system
from its Greek counterpart is the anchor point: where Hellenistic and later
Western astronomy measures longitude from the vernal equinox (a moving point
that drifts westward through the stars at roughly 50 arcseconds per year due to
the precession of Earth's rotation axis), Indian astronomy measures from fixed
stars.  A position is therefore \emph{sidereal}---\textit{nirayana}, ``without
the \textit{ayana}'' or solstitial shift---rather than tropical
(\textit{s\=ayana}).

The distinction is consequential.  Over the roughly two millennia since the
Indian and Hellenistic systems last agreed on the position of the vernal equinox
(approximately 285~CE by the most widely used \textit{Lahiri} \textit{ayanamsha},
though estimates range from 285 to 560~CE depending on the \textit{ayanamsha}),
precession has opened a gap of nearly 24° between tropical and sidereal
longitudes.  A planet at 10° tropical Aries sits near 16° sidereal Pisces.
Every translation between Indian and Chinese sources must contend with this
offset \cite{pingree1981jyotih}.

The twelve \textit{r\=asi}---\textit{Mesha} (Aries) through \textit{M\={\i}na}
(Pisces)---serve as the containers within which planetary positions are reported.
They are equal divisions, not constellations of variable extent, so the
\textit{r\=asi} framework is a clean mathematical abstraction.  Each
\textit{r\=asi} is further subdivided into \textit{nakshatra} (lunar mansion)
segments, linking the zodiacal and lunar-mansion systems in a unified
architecture.  It was this architecture---ecliptic plane, 360° circle,
equal-division signs, nine celestial bodies---that Tang dynasty translators would
import wholesale into China in 718~CE.

\subsection*{Rahu and Ketu: The Eclipse Demons}

Within the nine \textit{graha} (``seizers'' or celestial influences) of Indian
astronomy, two occupy a peculiar metaphysical status.  Rahu and Ketu are not
visible planets but mathematical points: the two nodes where the Moon's orbital
plane intersects the ecliptic.  At these intersections, and only at these
intersections, solar and lunar eclipses are geometrically possible.  Because
eclipses were among the most feared celestial events in ancient societies, the
nodes acquired a mythology commensurate with their power.

The origin story is preserved in the \textit{Mah\=abh\=arata} and the
\textit{Bh\=agavata Pur\=ana}.  During the churning of the cosmic ocean the gods
produced \textit{am\d{r}ta}, the nectar of immortality.  A demon named
Svarbh\=anu disguised himself among the gods and drank from the vessel.  The Sun
and Moon recognized the impostor and alerted Vishnu, who immediately severed
Svarbh\=anu's head from his body with his discus.  Because the demon had already
swallowed the nectar, however, both head and tail survived as immortal fragments:
the head became Rahu, the tail became Ketu.  Eternally enraged, they periodically
swallow the Sun and Moon---producing eclipses---only for the luminaries to emerge
from the severed neck moments later.

This mythological account maps precisely onto the astronomical reality.  Rahu,
the ascending (north) lunar node, is where the Moon crosses from south to north
of the ecliptic; Ketu, the descending (south) node, is where it crosses from
north to south.  The two nodes are always exactly 180° apart---head and tail of
the same orbital intersection---and they precess westward around the ecliptic
with a period of approximately 18.6 years.  The identification is astronomically
unambiguous and remained stable across Indian tradition for more than a millennium
\cite{pingree1981jyotih}.

One complication deserves note before this Indian baseline is left behind.  The
Sanskrit word \textit{ketu} originally meant ``banner,'' ``bright appearance,''
or ``radiant body,'' and in earlier Vedic and post-Vedic literature it referred
primarily to comets \cite[p.~47]{kotyk2018sinicization}.  The identification of
\textit{ketu} specifically with the descending lunar node is a later development,
consolidated in the \textit{siddh\=anta} tradition but never entirely displacing
the older, broader sense.  This semantic flexibility---a name that had already
migrated once, from comet to node---prefigures the second migration that would
occur in Tang dynasty China, when \zh{計都} moved from node to apogee.

\begin{table*}[ht]
\centering
\small
\caption{Rahu and Ketu: Indian original vs.\ mature Chinese identification.}
\label{tab:rahu-ketu}
\begin{tabular}{@{}lll@{}}
\toprule
Body & Indian (\textit{Jyotish}) & Chinese (\zh{果老星宗}) \\
\midrule
\zh{羅睺} (Rahu) & Ascending node & Ascending node \\
\zh{計都} (Ketu) & Descending node & Osculating lunar apogee \\
\zh{月孛} (Yuebei) & --- & Mean lunar apogee \\
\zh{紫氣} (Ziqi) & --- & Fictitious ($\sim$28-year cycle) \\
\bottomrule
\end{tabular}
\end{table*}

\subsection*{The Clarity That Makes Transformation Visible}

The Indian system's identification of Rahu and Ketu is admirably
unambiguous. Two nodes, always opposite, always moving at the same rate,
always causally connected to eclipses: the conceptual architecture is
self-consistent and empirically grounded.  Ketu is the descending node,
full stop.  There is no competing function for it to perform, no
ambiguity about its period or direction of motion.

Table~\ref{tab:rahu-ketu} places this Indian baseline alongside the
mature Chinese identification that will be traced through the following
sections.  The contrast is arresting.  In the Chinese \zh{七政四餘}
system, \zh{羅睺} (Rahu) retains its Indian identity as the ascending
node, but \zh{計都} has been detached from the descending node
altogether and reattached to the osculating lunar apogee---a point that
moves in the \emph{same} direction as the Moon rather than opposite to
it, has a period of roughly 8.85 years rather than 18.6 years, and has
no direct causal relationship to eclipses.  A new body, \zh{月孛}
(\textit{Yueb\`ei}, ``lunar comet''), has been invented to track the
mean apogee, and a fourth fictitious point, \zh{紫氣} (\textit{Ziq\`{\i}},
``purple vapour''), has been added with a $\sim$28-year cycle that has
no Indian precedent.

This proliferation did not happen overnight.  It was the cumulative
product of successive transformations---transliteration, reidentification,
indigenous restructuring, and hybridization---that will be traced
dynasty by dynasty, beginning with the first moment of contact in
718~CE.

% === §3 The Tang Transmission =================================================
\section{The Tang Transmission: 718--806~CE}\label{sec:tang}
% ~1,500 words — Task 7

\subsection*{Transformation Mode~1: Transliteration}

The first mode of transformation is the simplest: foreign astronomical
concepts are adopted with minimal modification.  Names are transliterated,
tables are copied, procedures are translated as literally as the target
language permits.  The receiving culture does not yet possess the
institutional confidence to reshape what it has received; it is still
in the posture of a student before a prestigious foreign tradition.
This is the characteristic mode of the early Tang reception of Indian
astronomy.

\subsection*{The \textit{Jiuzhi li} of 718~CE}

The pivotal document is the \textit{Jiuzhi li} \zh{九執曆}
(\textit{Navagraha-karana}, ``Calendar of the Nine Seizers''), translated
in 718~CE by Gautama Siddh\=artha (\zh{瞿曇悉達}), an astronomer of
Indian descent whose family had settled in Chang'an over the preceding
generation.  Gautama served in the Tang Bureau of Astronomy
(\zh{太史局}) and brought to his translation both fluency in Sanskrit
astronomical texts and direct access to the imperial court's demand for
accurate calendrical and astrological computation
\cite{yabuuchi1989sui,kotyk2018sinicization}.

The \textit{Jiuzhi li} introduced into China a package of Indian
astronomical innovations that had no precedent in the native tradition.
Most consequential for positional astronomy was the adoption of a
360-degree circle, replacing the Chinese convention of dividing the
ecliptic into 365.25 degrees (one degree per day of the solar year).
Alongside this came decimal notation for fractional positions, sine
tables computed at 3.75° intervals (a characteristic feature of the
Indian \textit{siddh\=anta} tradition that derives ultimately from
Hellenistic chord tables), and epicyclic planetary models for computing
the true positions of the five visible planets from their mean motions.

Among the nine \textit{graha} that give the text its name, two were
purely mathematical points: Rahu (\zh{羅睺}) and Ketu (\zh{計都}).
Their Chinese names are phonetic transliterations---\zh{羅睺} for
``Rahu'' and \zh{計都} for ``Ketu''---the clearest possible signal that
the Tang astronomers were importing the concept as a foreign package
rather than mapping it onto any pre-existing Chinese category.  Both
entered Chinese astronomical vocabulary simultaneously, both designated
lunar nodes, and both were assigned the same westward precessional
motion of approximately 3.06° per month
\cite{yabuuchi1989sui,kotyk2018sinicization}.

The \textit{Jiuzhi li} also introduced the ecliptic as an explicit
reference plane.  Chinese astronomers had for centuries measured
celestial positions in equatorial coordinates: right ascension expressed
as entry into one of the twenty-eight lunar mansions (\zh{入宿度},
\textit{r\`uxiud\`u}) and polar distance (\zh{去極度}, \textit{q\`uj\'ud\`u}).
The ecliptic was not ignored in Chinese astronomy---the twelve
\textit{ci} (\zh{次}) had long served as a solar-position scheme---but
it was not the primary framework for planetary positions.  The
\textit{Jiuzhi li} brought with it a fully ecliptic-centric positional
system, creating an immediate tension between the imported framework and
the native one that would not be resolved for nearly six centuries.

\subsection*{The Iranian Intermediary}

The Tang transmission was not, however, a clean direct line from India
to China.  Kotyk's detailed philological analysis demonstrates that the
\textit{Jiuzhi li} and the astrological texts that accompanied it show
unmistakable signs of Persian and Sasanian intermediary influence
\cite{kotyk2018sinicization,kotyk2017buddhist}.  Calculation procedures,
parameter values, and terminology in several passages align more closely
with Persian \textit{z\={\i}j} tradition than with the Indian
\textit{siddh\=anta} sources from which the system ultimately derives.

Mak's recent study of the Persian astronomical channel into Tang China
strengthens this picture considerably \cite{mak2024persian}.  The
Sasanian empire maintained sophisticated astronomical institutions, and
Persian court astronomers had been assimilating and adapting Indian
methods for at least two centuries before the Tang transmission.  When
Indian astronomical ideas arrived in Chang'an, they had already been
filtered through a Persian cultural lens that introduced modifications---
in nomenclature, in parameter selection, in the relative prominence
accorded different celestial bodies---that are not always easy to
disentangle from the original Indian substrate.

This Iranian intermediary matters for the history of Ketu's
transformation.  If modifications to the node--apogee identification
were already under way in Persian astrology, the Chinese reidentification
may reflect not a purely Chinese innovation but the endpoint of a
gradual drift that began somewhere along the Silk Road between India
and Chang'an.  The evidence is not yet conclusive, but the Persian
channel must remain in view as a possible locus of the first
reidentification.

\subsection*{The \zh{七曜攘災訣} and the Apogee Identification (806~CE)}

The second key Tang text is the \textit{Qiy\`ao r\`angz\=ai ju\'e}
\zh{七曜攘災訣} (``Formulae for Averting Disasters Caused by the Seven
Luminaries''), a Buddhist astrological manual compiled by a
Br\=ahma\d{n}a priest from western India sometime after 806~CE, the
epoch year (Year~1 of the Yuanhe \zh{元和} reign period) embedded in its
planetary tables.  The text would not survive in Chinese repositories;
it was brought to Japan in 865~CE by the Japanese monk Shuy\=u
(\zh{宗叡}) and preserved within the \textit{Sukuy\=od\=o}
(\zh{宿曜道}, ``Way of the Lunar Mansions and Stars'') tradition, where
it remained effectively unknown to Chinese and Western scholarship until
modern times \cite{kotyk2017buddhist}.

The \textit{Qiy\=ao r\`angz\=ai ju\'e} contains multi-decade ephemerides for
all nine \textit{graha}, including a Rahu table spanning 93~years and a
Ketu table spanning 62~years.  In 1994, the Chinese historian of science
Niu Weixing subjected the Ketu ephemeris to a systematic computational
reanalysis, comparing its tabulated positions against two candidates:
the descending lunar node as defined in Indian \textit{siddh\=anta}
astronomy, and the osculating lunar apogee \cite{niu1995ketu}.  The
result was unambiguous: the Ketu ephemeris fits the motion of the
\emph{apogee}, not the node.  The period, the direction of motion, and
the epoch position all correspond to the lunar apogee; none correspond
to the descending node.

This finding establishes a terminus ante quem for the reidentification.
By 806~CE---within roughly one century of the \textit{Jiuzhi li}
translation---\zh{計都} had already ceased to designate the descending
lunar node and had been reassigned to track the osculating apogee.  The
name persisted; the referent changed completely.

\subsection*{Coordinate Thread: The Ecliptic in an Equatorial World}

The coordinate dimension of the Tang transmission is a story of
institutional coexistence rather than replacement.  The \textit{Jiuzhi li}
introduced an ecliptic framework, but the Bureau of Astronomy continued
to produce the official calendar (\textit{li} \zh{曆}) in the native
equatorial system.  Indian ecliptic methods were used for astrological
purposes---planetary hour lords, eclipse prediction, the \textit{graha}
positions in horoscopes---while the equatorial system retained its
authority for the state rituals tied to the solar terms and the
twenty-eight lunar mansions.

The result was a bifurcated institutional landscape.  Buddhist
astrologers working from Indian and Persian sources operated in ecliptic
coordinates; the official Bureau computed in equatorial ones.  The two
systems were not formally reconciled, and conversion procedures between
them were approximate at best.  This bifurcation, established in the
early eighth century, would persist as an unresolved tension through the
Song dynasty, when the first serious attempt at a unified ecliptic
sidereal system for Chinese \textit{qizheng siyu} (\zh{七政四餘},
Seven Governors and Four Remainders) astrology would finally be made.

\subsection*{Identity Thread: Mode~2 Already in Motion}

The \textit{Qiy\=ao r\`angz\=ai ju\'e}'s apogee ephemeris means that by
the end of the Tang period, both transformation modes are already
visible.  Mode~1 (transliteration) is represented by the \textit{Jiuzhi
li}: foreign names and methods are adopted wholesale, with \zh{羅睺}
and \zh{計都} entering Chinese as phonetic loan words for Indian
concepts.  Mode~2 (reidentification) is already under way: the name
\zh{計都} has been retained but silently detached from its original
referent.

What is striking about this early reidentification is its lack of
documentary justification.  No Tang text known to survive explains
\emph{why} Ketu should track the apogee rather than the node.  The
change is not announced, not argued for, not defended.  It appears in
the ephemeris tables as a fait accompli, suggesting that it may have
occurred upstream---in the Persian intermediary tradition, or in a
lost stratum of Indian regional astronomy---rather than being a
conscious Chinese innovation.  Whatever its origin, the reidentification
was consequential: it would shape every subsequent Chinese treatment of
the Four Remainders, creating the distinctive asymmetry---Rahu as node,
Ketu as apogee---that defines the mature \zh{七政四餘} system and
persists in software to this day.

% === §4 The Song System =======================================================
\section{Ecliptic Sidereal: The Song System (\zh{三辰通載})}
\label{sec:song}
% ~1,200 words — Task 8

\subsection*{Transformation Mode~2: The Post-Transliteration State}

The Tang transmission planted Indian astronomical seeds in Chinese soil.
The Song dynasty (960--1279~CE) represents the first flowering: a codified,
institutionalized system that preserved the imported framework with remarkable
fidelity.  Before the indigenous restructuring of the Yuan and the Jesuit
correction of the Qing, there existed a Chinese \zh{七政四餘} practice that
looked, structurally, very much like what a medieval Indian astrologer would
recognize.  The \zh{三辰通載} (\textit{S\={a}nch\'{e}n t\={o}ngz\`ai},
``Comprehensive Record of the Three Luminaries'') is our primary evidence for
this baseline state.

\subsection*{The \zh{三辰通載} as Earliest Codification}

The \zh{三辰通載} is the earliest extant Chinese text to codify the full
\zh{七政四餘} system in a form recognisably continuous with its Indian
antecedents.\footnote{Direct access to the \zh{三辰通載} text remains limited;
our characterization relies on secondary analyses, principally
\citealt{ho2003astrology} and \citealt{mak2014pseudoplanets}.}
Attributed to the Song period and concerned with what Chinese tradition called
the ``three luminaries''---Sun, Moon, and stars---the text assembles the
seven visible planets (\zh{七政}: Sun, Moon, Mercury, Venus, Mars, Jupiter,
Saturn) together with the four ``remainders'' (\zh{四餘}) into a unified
predictive scheme.

What makes the \zh{三辰通載} a baseline rather than merely an early text
is precisely what it does \emph{not} do.  It does not translate the imported
ecliptic frame into native equatorial coordinates.  It does not recast
planetary dignities in terms of the Five Phases (\zh{五行}).  It does not
subject the pseudo-planets to a programme of sinicization.  The system
is cosmopolitan in the most literal sense: it operates within coordinate
conventions and interpretive categories that arrived from the west and
were not yet remade into something distinctively Chinese.

\subsection*{Coordinate Thread: The Ecliptic Sidereal Frame}

The \zh{三辰通載} employs what we designate the \zh{黃道恆星製}
(\textit{hu\'{a}ngd\`ao h\'{e}ngx\=ing zh\`{\i}}, ``ecliptic sidereal
system'').  Planetary longitudes are measured along the ecliptic from a
fixed sidereal zero-point---the same fundamental choice that underpins
Sanskrit \textit{siddh\={a}nta} astronomy.  The twelve zodiacal houses
(\zh{十二宮}) are equal 30\textdegree{} divisions of the ecliptic, each
mapped to one of the twelve Earthly Branches (\zh{地支}): \zh{子} (Zi)
for Aquarius, \zh{丑} (Chou) for Capricorn, and so on in the familiar
sequence.  This direct correspondence between the Indian zodiac and the
Chinese duodenary cycle is not a coincidence; it is a deliberate
translation, executed during the Tang transmission, that anchored the
foreign framework to an existing Chinese indexing system
\citep{kotyk2018sinicization}.

Within each house, the interpretive machinery is recognisably Jyotish.
Aspects (\zh{刑沖合會}) correspond to the Indian \textit{dr\d{s}\d{t}i}
doctrine.  Planetary dignity---exaltation, detriment, and the like---maps
to \textit{ucca} and \textit{n\={\i}ca} as transmitted through Persian
intermediaries.  Crucially, no Five-Phase (\zh{五行}) overlay is present.
The planets are not assigned to Wood, Fire, Earth, Metal, and Water in
their astrological capacities; they remain the seven Greco-Indian wanderers,
differentiated by their behavior within the ecliptic frame rather than by
their membership in a Chinese cosmological scheme.

Two structural features deserve special notice.  First, the text employs
the \zh{七曜值日} (``Seven Luminaries Day-Rule''), the planetary week in
which each day and hour is governed in turn by one of the seven planets.
This is the Indian \textit{v\={a}ra} (planetary weekday) doctrine, which
reached China via the same Sogdian-Persian channel that transmitted the
\textit{Jiuzhi li} \citep{mak2024persian}.  The seven-day week is not
a native Chinese institution; its presence in the \zh{三辰通載} is a
direct inheritance from the Tang transmission, unmodified.

Second, the text contains the planetary period system known as
\zh{洞微百六限} (\textit{D\`ongw\=ei b\`aili\`u xi\`an},
``Penetrating-Subtlety Hundred-and-Six Limits'').  This is the Chinese
adaptation of the Indian \textit{da\'{s}\={a}} (time-lord) system, in which
each planet governs a defined span of a nativity's life in sequence.
The specific allocation of years and the sequencing of the seven governors
in the \zh{洞微百六限} shows closer affinity with the Indian
\textit{Yavanaj\={a}taka} tradition than with later purely Chinese
alternatives \citep{ho2003astrology}.

\subsection*{Identity Thread: The Four Remainders Crystallize}

The most consequential development in the \zh{三辰通載} for the purposes
of this paper is the crystallization of the \zh{四餘} as a canonical set.
By the Song period the four fictitious bodies have settled into their
classical configuration: \zh{羅睺} (Rahu, ascending node), \zh{計都}
(Ketu), \zh{月孛} (\textit{Yu\`eb\`ei}, ``Lunar Comet,'' mean apogee),
and \zh{紫氣} (\textit{Ziq\`{\i}}, ``Purple Vapour,'' a 28-year
intercalary point).

The identity of \zh{計都} in this period is significant.  As noted in
Section~\ref{sec:tang}, Niu Weixing's analysis of the 806~CE ephemeris
establishes that the apogee reidentification was already in place by
the end of the Tang.  The Song system inherits this identification
and embeds it in the stable four-body set.  The doubled apogee
tension---\zh{計都} as osculating apogee and \zh{月孛} as mean
apogee simultaneously present in the same system---is therefore a
structural feature of \zh{七政四餘} from its earliest codification,
not a later innovation.  We return to this tension in
Section~\ref{sec:yuan}.

\subsection*{Significance: A Baseline for Measuring Transformation}

The \zh{三辰通載} system is significant precisely as a \emph{baseline}.
It preserves the maximum Indian DNA that the Chinese tradition will
ever retain simultaneously in a single codified text.  Ecliptic
coordinates, sidereal zero-point, Jyotish-style house structure,
\textit{v\={a}ra} day-rule, \textit{da\'{s}\={a}}-derived period
system---all are present, with only modest adaptation.

This makes the \zh{三辰通載} methodologically invaluable.  The
Yuan equatorial reform (Section~\ref{sec:yuan}), the Ming quasi-ecliptic
synthesis (Section~\ref{sec:ming}), and the Jesuit tropical correction
(Section~\ref{sec:jesuit}) will successively overwrite exactly these
features.  Each subsequent transformation can be precisely characterized
by reference to what it departs from.  Without the Song baseline, the
later changes would be difficult to measure; with it, the direction and
magnitude of each transformation become legible.

In the language of the six-mode framework introduced in
Section~\ref{sec:intro}, the \zh{三辰通載} represents the terminal
state of Mode~1 (transliteration) and the opening condition for
Mode~3 (indigenous restructuring) that will follow in the Yuan.

% === §5 The Yuan Reform =======================================================
\section{Equatorial Sidereal: The Yuan Reform (\zh{鄭氏星案})}
\label{sec:yuan}
% ~1,500 words — Task 8

\subsection*{Transformation Mode~3: Indigenous Restructuring}

If Mode~1 (transliteration) is assimilation without digestion, and
Mode~2 establishes the baseline that transliteration produces, then
Mode~3 is the moment when the host culture begins actively reshaping
what it has received.  Indigenous restructuring does not reject the
imported content; it preserves the pseudo-planets, the period system,
the interpretive categories.  What it replaces is the \emph{form}---the
coordinate system in which that content is expressed.  The Yuan dynasty
(1271--1368~CE) provides the most dramatic instance of this mode in
the entire history of Chinese \zh{七政四餘}.

\subsection*{Historical Context: Guo Shoujing and the \zh{授時曆}}

The Yuan was an era of exceptional astronomical ambition.  Kublai
Khan's court sponsored simultaneous engagement with both the indigenous
Chinese tradition and the Islamic astronomical science arriving via
the westward connections of the Mongol empire.  An Islamic Astronomical
Bureau (\zh{回回司天監}, \textit{Huihu\'{i} S\={i}ti\={a}nj\={a}n})
operated in parallel with the traditional Chinese bureau, bringing
\textit{z\={\i}j} tables and armillary instruments of Persian and
Arabic origin into direct institutional contact with Chinese
practice \citep{mak2024persian}.

Within this environment, the astronomer Guo Shoujing (\zh{郭守敬},
1231--1316) led the most technically ambitious calendar reform in
Chinese history.  The resulting \zh{授時曆} (\textit{Sh\`oush\`{\i}
l\`{\i}}, ``Season-Granting Calendar,'' 1280) achieved solar-year
accuracy not surpassed in China for three centuries
\citep{sivin2009granting}.  Its technical foundation was the
equatorial coordinate system: Guo Shoujing designed a suite of
precision instruments, most famously the \zh{簡儀} (\textit{ji\v{a}n
y\'{i}}, ``simplified armillary''), that measured celestial positions
directly in right ascension and declination rather than converting
through the ecliptic.

The \zh{授時曆} also solved, for the first time in Chinese astronomy,
the mathematical problem of converting between ecliptic and equatorial
coordinates with precision.  Wagner's reconstruction of Guo Shoujing's
conversion algorithm demonstrates that the \zh{授時曆} employed a
spherical-trigonometric method---unprecedented in Chinese mathematical
astronomy---to map positions between the two frames
\citep{wagner2024gsj}.  Whether this mathematical capacity directly
influenced the \zh{七政四餘} practitioners who produced the
\zh{鄭氏星案} remains an open question; but the institutional and
intellectual environment in which equatorial measurement became the
dominant mode is the necessary precondition for the coordinate shift
we document.

\subsection*{The Text: \zh{鄭氏星案}}

The \zh{鄭氏星案} (\textit{Zh\`engsh\`{\i} x\={\i}ng'\`an},
``Zheng's Stellar Cases'') is a Yuan-dynasty \zh{七政四餘} manual
that represents the first systematic application of equatorial
coordinates to the inherited Indian astrological framework.\footnote{The
equatorial sidereal characterization of the \zh{鄭氏星案} awaits full
verification against the primary text; our characterization of its
coordinate system draws on secondary analyses, notably
\citealt{ho2003astrology} and \citealt{niu2010fourremains}.}
The title ``Stellar Cases'' suggests a practitioner's casebook of
worked examples---a format that carries the reader step by step through
astrological delineation---rather than a theoretical treatise.  This
makes it all the more striking that so fundamental a feature as the
coordinate system was reformulated.

\subsection*{Coordinate Thread: From Ecliptic to Equatorial}

The \zh{鄭氏星案} employs the \zh{赤道恆星製}
(\textit{ch\`{\i}d\`ao h\'{e}ngx\=ing zh\`{\i}},
``equatorial sidereal system'').  Where the Song \zh{三辰通載}
measured planetary longitudes along the ecliptic, the Yuan text
measures positions in right ascension, anchored to the native
Chinese system of twenty-eight lunar mansions (\zh{二十八宿},
\textit{\`ersh\'{i}b\={a} xi\`u}).  The twelve zodiacal houses,
in the Song system clean 30\textdegree{} ecliptic arcs,
become divisions of the celestial equator, irregular in ecliptic
terms but aligned with the observational infrastructure that Chinese
astronomy had built over millennia.

The conceptual consequences are far-reaching.  The ecliptic is the
plane of the zodiac; the zodiacal constellations that define Aries,
Taurus, Gemini, and the rest are ecliptic phenomena.  Once the
coordinate backbone is moved to the equator, the zodiacal constellation
concept does not merely become imprecise---it becomes inapplicable.
The \zh{鄭氏星案} does not attempt to preserve it.  What survives the
coordinate transformation is the functional infrastructure of
prediction: the seven governors, the four remainders, the period
system, the aspects.  What does not survive is the ecliptic framework
that gave the zodiacal houses their geometric meaning.

In the terminology of our six-mode framework, this is Mode~3 in its
purest form.  The \emph{content} of the imported system---pseudo-planets,
time-lords, interpretive techniques---persists intact.  The
\emph{form}---the coordinate operating system in which that content
runs---is replaced by the native Chinese alternative.  An analogy
from the history of computing is instructive: astronomical
``data'' has been successfully ported to a new ``operating system,''
but applications written natively for the old platform no longer
run without modification.  The zodiacal house system, native to the
ecliptic, cannot run on the equatorial platform without
fundamental redesign.

\subsection*{The Mechanism of Conversion}

How were positions translated from the old ecliptic sidereal frame
to the new equatorial sidereal frame in practice?  The \zh{授時曆}'s
conversion algorithm, reconstructed by Wagner, provides one plausible
pathway \citep{wagner2024gsj}.  Guo Shoujing's method works forward
from equatorial observations to ecliptic positions for calendar
purposes; the inverse problem---converting existing ecliptic-referenced
\zh{七政四餘} tables into equatorial coordinates---is mathematically
equivalent.

Whether \zh{七政四餘} practitioners had direct access to the
\zh{授時曆}'s computational machinery, or worked from
approximately-converted tables, or simply re-observed planetary
positions against the lunar mansion system, cannot be established
from the surviving text with certainty.  What can be established
is the outcome: the Yuan \zh{七政四餘} practitioner no longer
worked in ecliptic longitude but in a right-ascension equivalent
measured against the 28 lodges.  The practical effect was to make
\zh{七政四餘} positions directly comparable with, and perhaps
checkable against, the standard observational output of the
official astronomical bureau.

\subsection*{Identity Thread: The Doubled Apogee Persists}

The Yuan reform touches coordinates but not identity.  The \zh{四餘}
continue within the equatorial framework with their Song-period
identifications undisturbed.  \zh{羅睺} remains the ascending
node.  \zh{計都} remains the osculating lunar apogee.  \zh{月孛}
(\textit{Yu\`eb\`ei}) continues to track the mean apogee.
\zh{紫氣} (\textit{Ziq\`{\i}}) continues its 28-year cycle.

The doubled-apogee tension noted in Section~\ref{sec:song}---two
bodies, \zh{計都} and \zh{月孛}, both apsidal, one osculating and
one mean---therefore persists through the coordinate transformation.
This is not a logical contradiction within the system's own terms;
osculating and mean apogee are genuinely distinct quantities that
diverge by up to several degrees at any given moment.  But
from the perspective of an observer comparing the system with
Indian \textit{Jyotish}, the presence of two apogees where the
Indian system has one (the descending node) and one apogee
(\textit{mand\={o}cca}) requires explanation.  The Yuan text,
like the Song text before it, does not supply one.  The doubled
apogee is a structural feature that will be inherited, unchanged
in its identity if not in its coordinates, by the Ming synthesis.

\subsection*{Significance: The Most Dramatic Transformation}

Of the four coordinate regime changes traced in this paper, the
Yuan shift from ecliptic to equatorial is the most dramatic in
geometric terms.  The ecliptic and the equator intersect at an
angle of approximately 23.5\textdegree{} (the obliquity of the
ecliptic), which means that a position expressed in ecliptic
longitude and latitude cannot be converted to equatorial right
ascension and declination by a simple additive offset: the
conversion requires a rotation of the coordinate frame.  The
Song-to-Yuan transition therefore involved a genuine geometric
transformation of the entire positional database, not merely a
recalibration of a zero-point.

What makes this transformation analytically tractable is precisely
the stability of the content layer.  Because the pseudo-planets,
interpretive rules, and period allocations were preserved unchanged
across the coordinate shift, the transformation is detectable in
retrospect: one can compare the positions predicted by the Song
ecliptic model with those predicted by the Yuan equatorial model
for the same historical date, and the systematic offset will reveal
the nature of the conversion actually performed.  The computational
verification in Section~\ref{sec:computation} exploits exactly
this property.

The Yuan reform also opens a question that the text itself does
not answer: was the equatorial reframing driven by intellectual
preference---a desire to align \zh{七政四餘} with the prestige
of Guo Shoujing's new astronomical infrastructure---or by
practical necessity, as practitioners trained in the equatorial
tradition found ecliptic-referenced tables difficult to use?
The answer likely involves both, and the ambiguity is itself
historically significant.  What is clear is that after the Yuan
reform, the ecliptic sidereal frame of the Song baseline was
gone, replaced by a coordinate system whose deepest roots were
not Indian or Persian but natively Chinese.

% === §6 The Ming Synthesis ====================================================
\section{Quasi-Ecliptic Sidereal: The Ming Synthesis (\zh{果老星宗})}
\label{sec:ming}
% ~1,800 words — Task 9

\subsection*{Transformation Mode~4: Hybridization}

The Ming dynasty did not merely inherit the Yuan equatorial frame.  It
transformed it a second time, producing a coordinate system with no
counterpart in any other astronomical tradition---a system that looks
ecliptic-like on its surface but is derived from equatorial
measurements underneath.  This hybridization is the paper's most
important case study, and it is the juncture at which the
\zh{七政四餘} tradition achieves maximum distance from its Indian
original in both coordinate conventions and interpretive content.

\subsection*{The Canonical Texts}

The primary vehicle for this transformation is the
\zh{果老星宗} (\textit{Gu\v{o}l\v{a}o x\={\i}ngz\={o}ng}, ``Star-School
of the Elder of the Fruit''), the Ming-dynasty canonical text of the
\zh{七政四餘} tradition.\footnote{The \zh{果老星宗} takes its name
from Zhang Guolao (\zh{張果老}), one of the Eight Immortals of Daoist
legend and a figure credited, in mythologised form, with transmitting
astral knowledge from antiquity.  The attribution is hagiographic
rather than historical.}  The text compiles computational rules,
interpretive glosses, and worked examples for all eleven bodies of
the mature system.  It is the text that most modern Chinese
astrologers regard as the authoritative statement of ``standard''
\zh{七政四餘} practice, and its coordinate conventions are therefore
the ones most consequential for the living tradition.

A near-contemporary source of comparable authority is the
\zh{星學大成} (\textit{X\={\i}ngxu\'e d\`ach\'eng}, ``Great Compendium
of Astral Studies'') by Wan Minying \zh{萬民英} (1521--1603)
\cite{ho2003astrology}.  Wan's compilation is more systematic in its
presentation of computational methods, and the two texts are
frequently cited together as the twin pillars of Ming
\zh{七政四餘} scholarship.  Where they agree---and on the coordinate
conventions they agree---we can speak of a Ming consensus.

\subsection*{Coordinate Thread: Defining the \zh{似黃道恆星製}}

The coordinate system of the Ming texts is the
\zh{似黃道恆星製} (\textit{s\`{\i} hu\'angd\`ao h\'engx\={\i}ng
zh\`{\i}}, ``quasi-ecliptic sidereal system'').  The name encodes
both an assertion and a qualification: the system is \textit{like}
(\zh{似}) the ecliptic---it maps positions onto zodiac-style signs
and houses---but the underlying measurements are taken along the
equator, not the ecliptic.  No other astronomical tradition has
produced such a hybrid.

\subsubsection*{Mathematical Definition}

The construction proceeds in three steps.  First, the equator is
divided into twelve equal houses of $30\degree$ each, beginning at
the vernal equinox.  Second, each equatorial boundary is projected
onto the ecliptic plane using the standard transformation between
right ascension $\alpha$ and ecliptic longitude $\lambda$ at zero
ecliptic latitude ($\beta = 0$):
\begin{equation}
  \tan\alpha = \cos\varepsilon \cdot \tan\lambda
  \label{eq:ra-to-ecl}
\end{equation}
where $\varepsilon \approx 23.5\degree$ is the obliquity of the
ecliptic.  Inverting to obtain the ecliptic boundary corresponding to
each equatorial house boundary:
\begin{equation}
  \lambda = \arctan\!\left(\frac{\sin\alpha}{\cos\alpha \cdot \cos\varepsilon}\right)
  \label{eq:ecl-boundary}
\end{equation}
Third, planetary positions, originally computed in equatorial
coordinates, are re-expressed in terms of whichever ecliptic house
they fall into under the projected boundaries.

The result is that the twelve ``signs'' look ecliptic---they are
labelled with the same zodiacal names as in any ecliptic
system---but their boundaries are not equally spaced in ecliptic
longitude.  The unequal spacing is a direct consequence of the
obliquity $\varepsilon$.  Houses near the equinoxes, where the
equator and ecliptic are most nearly parallel in their rate of
longitude change, acquire wider ecliptic extents; houses near the
solstices, where the relationship is most compressed, acquire
narrower extents.

\subsubsection*{Computed Boundaries and the Symmetry Pattern}

Table~\ref{tab:quasi-ecliptic} gives the computed ecliptic boundaries
for all twelve houses under obliquity $\varepsilon = 23.5\degree$
(appropriate for approximately 1400~CE).

\begin{table}[ht]
\centering
\caption{Quasi-ecliptic (\zh{似黃道}) house boundaries compared to equal ecliptic divisions. Obliquity $\varepsilon = 23.5\degree$ (circa 1400~CE). Houses near equinoxes are wider ($\sim$32.2\degree); houses near solstices are narrower ($\sim$27.9\degree).}
\label{tab:quasi-ecliptic}
\small
\begin{tabular}{@{}rrrr@{}}
\toprule
House & RA boundary & Ecliptic boundary & Width \\
\midrule
1 & 0.0\degree & 0.00\degree & 32.19\degree \\
2 & 30.0\degree & 32.19\degree & 29.91\degree \\
3 & 60.0\degree & 62.10\degree & 27.90\degree \\
4 & 90.0\degree & 90.00\degree & 27.90\degree \\
5 & 120.0\degree & 117.90\degree & 29.91\degree \\
6 & 150.0\degree & 147.81\degree & 32.19\degree \\
7 & 180.0\degree & 180.00\degree & 32.19\degree \\
8 & 210.0\degree & 212.19\degree & 29.91\degree \\
9 & 240.0\degree & 242.10\degree & 27.90\degree \\
10 & 270.0\degree & 270.00\degree & 27.90\degree \\
11 & 300.0\degree & 297.90\degree & 29.91\degree \\
12 & 330.0\degree & 327.81\degree & 32.19\degree \\
\bottomrule
\end{tabular}
\end{table}

The symmetry pattern is instructive.  Houses~1, 6, 7, and~12 (near
the equinoxes) are the widest at $\approx 32.19\degree$; houses~3,
4, 9, and~10 (near the solstices) are the narrowest at
$\approx 27.90\degree$.  The intermediate houses (2, 5, 8, 11) fall
at $\approx 29.91\degree$.  The pattern has period-two symmetry
because the obliquity effect is symmetric about both equinox lines:
the ecliptic and equator cross at the same angle at Aries~$0\degree$
and at Libra~$0\degree$, and the same compression applies at the two
solsticial points.

The maximum deviation from the equal $30\degree$ division is
$\approx \pm 2.19\degree$.  This seems small.  It is not
astronomically negligible.  A body at the boundary between two houses
can be placed in either house depending solely on which convention is
applied.  For slow-moving bodies such as Saturn (average motion
$\approx 0.033\degree$/day), the $\pm 2.19\degree$ discrepancy
corresponds to a positional ambiguity of roughly 66 days.  For the
nodes (average regression $\approx 0.053\degree$/day) the ambiguity
is approximately 41 days.  The quasi-ecliptic frame therefore
introduces a systematic, predictable source of inter-system
discrepancy that is distinct from---and additive with---the
coordinate-zero-point discrepancies discussed in
Sections~\ref{sec:song} and~\ref{sec:jesuit}.

\subsubsection*{Why the System Is Novel}

It is worth pausing on this point.  Indian astrology uses ecliptic
coordinates directly.  Greek astrology uses ecliptic coordinates
directly.  Persian and Arabic astrology---the traditions that served
as intermediaries for the Tang transmission---use ecliptic
coordinates directly.  The Yuan reform replaced ecliptic with
equatorial coordinates.  The Ming synthesis did something none of
these did: it used equatorial measurements to generate apparently
ecliptic house divisions.  The \zh{似黃道恆星製} is a genuinely new
coordinate species, not a corruption or a simplification of any
predecessor system.

\subsection*{Identity Thread: The Mature \zh{四餘} System}

The Ming texts present the most fully elaborated version of the
four-remainder (\zh{四餘}) system.  The four bodies are now formally
defined, their computational methods codified, and their astrological
significances established.  The two bodies whose identities were
contested in earlier periods now coexist in a stable, if
theoretically uneasy, configuration.

\zh{羅睺} (\textit{R\'ahu}) retains the identity it has held since
the Tang transmission: the ascending lunar node.  No reidentification
touches it.  It remains the closest of the four bodies to its Indian
original in both name and astronomical definition.

\zh{計都} (\textit{K\`et\`u}) is the osculating lunar apogee---the
reidentification already in place by the end of the Tang and
confirmed by the Song codification.  In the Ming texts this identity
is simply assumed, not argued; the apogee assignment has become the
tradition.

\subsubsection*{The Mean Apogee: \zh{月孛}}

\zh{月孛} (\textit{Yu\`eb\`o\`{\i}}, ``lunar comet'' or ``lunar
apoapsis'') is the mean lunar apogee---the apogee averaged over the
Moon's short-period perturbations and free of the oscillations that
distinguish the osculating apogee.  Its identification with Zhao
Youqin (\zh{趙友欽}, fl.\ late 13th--early 14th century) represents
one of the genuinely indigenous contributions to the system
\cite{niu2010fourremains}.  Zhao's derivation proceeded from careful
observation of the Moon's varying daily motion: slower motion near
apogee, faster near perigee.  By tracking the Moon's speed across
many lunations, the mean apogee longitude can be extracted without
requiring explicit knowledge of the underlying orbital geometry.
The result is a quantity distinct from \zh{計都} precisely because
it is mean rather than osculating.

\subsubsection*{The Purple Qi: \zh{紫氣}}

\zh{紫氣} (\textit{Z\v{\i}q\`{\i}}, ``Purple Qi'') is a fictitious
point with a cycle of approximately 28 years.  It has no
astronomical counterpart---it does not correspond to any orbital
element of the Moon, the Sun, or any planet---and it has no Indian
precedent.  It is a purely Chinese invention, its cycle apparently
chosen to be commensurate with Jupiter's 12-year cycle and Saturn's
29.5-year cycle in a numerologically satisfying way.  Unlike the
other three remainders, \zh{紫氣} cannot be confirmed or disconfirmed
by independent astronomical observation; its position is determined
entirely by the calendrical formula that defines it.

\subsubsection*{The Doubled Apogee Problem}

The Ming system thus contains two bodies that both track the lunar
apse line: \zh{計都} (osculating apogee) and \zh{月孛} (mean
apogee).  They inhabit the same orbital degree of freedom.  At any
given moment they differ by a quantity that is the difference between
the osculating and mean lunar apogee---a difference that varies over
the anomalistic month and reaches a maximum of several degrees.

This structural redundancy has never been satisfactorily resolved
within the tradition.  No Ming or subsequent text provides a rigorous
explanation of why both bodies are necessary, or what astrological
distinction should be drawn between them.  The doubled apogee
remains an open theoretical question \cite{niu2010fourremains}.

\subsubsection*{The Beginning of Bazi-ification}

One further development of the Ming period deserves mention, even
though its full consequences belong to a later section.  Mak (2014)
documents the emergence of a practice he terms ``bazi-ification'':
the derivation of \zh{四餘} positions from \zh{八字}
(\textit{b\={a}z\`{\i}}, ``eight characters'' or ``four pillars'')
numerology rather than from astronomical ephemerides
\cite{mak2014pseudoplanets}.  In this approach, a practitioner
determines the positions of the four remainders not by computing
their actual positions in the sky but by applying formulas that take
the birth date and time as input and yield positions by arithmetic
rule.  The positions so obtained are not guaranteed to match the
astronomical bodies they nominally represent.

This is the first step toward the algorithmic fossilization that
will be examined in Section~\ref{sec:fossil}.  At the Ming stage it
is a minority practice competing with genuine ephemeris-based
computation.  Its significance is not yet apparent.

\subsection*{Significance: High-Water Mark of Sinicization}

The Ming synthesis represents the maximum distance the
\zh{七政四餘} system ever achieved from its Indian original,
measured along both the coordinate and identity threads simultaneously.

On the coordinate thread, the \zh{似黃道恆星製} is the culmination
of a three-stage transformation.  The Song system preserved the
Indian ecliptic frame.  The Yuan reform replaced it with the native
Chinese equatorial frame.  The Ming synthesis performed a further
operation: it used the equatorial frame to construct an
\textit{apparent} ecliptic frame whose divisions are unequal in
ecliptic longitude.  The result superficially resembles the ecliptic
systems of the Indian, Greek, and Arabic traditions but is not
reducible to any of them.  Mode~4 (hybridization) is the appropriate
designation: the imported ecliptic appearance and the native
equatorial foundation are fused into something genuinely new.

On the identity thread, all four remainders are now present, two of
them (\zh{月孛} and \zh{紫氣}) without Indian precedent.  The
tradition has added bodies from native Chinese sources and retained
them in permanent parallel with the bodies inherited from India.

Yet the system has not yet fossilized.  The \zh{果老星宗} and
\zh{星學大成} both give computational procedures that, however
hybrid, are in principle checkable against the sky.  Positions are
derived from genuine orbital elements, even if the coordinate frame
in which they are expressed is novel.  The astronomical content
survives.  The next disruption will come not from Chinese synthesis
but from European imposition \cite{mak2014pseudoplanets,niu2010fourremains}.

% === §7 The Jesuit Disruption =================================================
\section{Ecliptic Tropical: The Jesuit Disruption (1645)}
\label{sec:jesuit}
% ~2,000 words — Task 9

\subsection*{Transformation Mode~5: Foreign Correction}

In 1645 Johann Adam Schall von Bell (\zh{湯若望}), a German Jesuit
astronomer serving as director of the imperial Bureau of Astronomy,
promulgated a reformed astronomical system under the new Qing dynasty.
The reform was the most violent single disruption in the history of
the \zh{七政四餘} tradition---more disruptive, in the short term,
than even the Tang transmission---because it changed both the
coordinate thread and the identity thread simultaneously, imposed
from outside the tradition by institutional authority rather than
emerging from internal Chinese scholarly development.

This is Mode~5: foreign correction.  Unlike Mode~3 (indigenous
restructuring, which was organic cultural pressure adapting an
imported system to native categories), Mode~5 is a deliberate,
externally authorized reset.  The new system was not proposed for
scholarly debate; it was mandated by imperial decree.

\subsection*{The Reform and Its Context}

The institutional context matters.  The late Ming dynasty had
suffered a series of calendar prediction failures that undermined
the credibility of the existing astronomical administration.  When
the Qing forces captured Beijing in 1644, Schall---who had been in
China since 1622 and had already demonstrated the accuracy of
European methods in predicting eclipses---was positioned to offer an
alternative.  In 1645 the Shunzhi Emperor promulgated the
\zh{時憲曆} (\textit{Sh\`{\i}xi\`an l\`{\i}}, ``Calendar of Edicts
Conforming to Heaven''), the new official calendar based on European
methods \cite{liu2020missionary}.

The computational infrastructure underlying the calendar was drawn
from the Tychonic system: Tycho Brahe's geo-heliocentric model,
which placed the planets in orbit around the Sun while maintaining
the Sun in orbit around the Earth.  From a positional astronomy
standpoint, the Tychonic system is mathematically equivalent to the
Copernican heliocentric model for the purposes of computing planetary
longitudes, and its predictions were substantially more accurate
than those of the Chinese \zh{授時曆} over the century of positional
drift that had accumulated since the Yuan reform \cite{jami2015calendar}.

For the \zh{七政四餘} system specifically, the reform imposed a
coordinate change of fundamental character: from sidereal to tropical.
The zero point of longitude shifted from fixed stars to the vernal
equinox.  The new system can be designated the
\zh{黃道迴歸製} (\textit{hu\'angd\`ao hu\'{\i}gu\={\i} zh\`{\i}},
``ecliptic tropical system'').

\subsection*{Coordinate Thread: The Ironic Completion}

The coordinate history of \zh{七政四餘} traces a closed trajectory
whose ironic structure is only visible from the vantage point of
the Qing reform.

The Song system used \textit{ecliptic sidereal} coordinates: the
ecliptic plane, zero point at a fixed star.  The Yuan reform
shifted to \textit{equatorial sidereal} coordinates: the equatorial
plane, zero point still at a fixed star.  The Ming synthesis
produced the \textit{quasi-ecliptic sidereal} hybrid: equatorial
measurements mapped to ecliptic-like divisions, zero point at a
fixed star.  The Qing reform imposed \textit{ecliptic tropical}
coordinates: the ecliptic plane, zero point at the vernal equinox.

The trajectory thus moves from ecliptic sidereal (Song) to equatorial
sidereal (Yuan) to quasi-ecliptic sidereal (Ming) and then snaps
back to the ecliptic plane (Qing)---but with the zero point
replaced.  The system ended in the ``opposite corner'' from where
it began: same plane as the Indian original (ecliptic), same plane
as the Song codification, but a different zero point.

The magnitude of the coordinate zero-point discrepancy is
substantial.  The tropical-sidereal difference is the precession
accumulated since the epoch at which the sidereal zero was defined.
By 1645, this amounts to approximately $24\degree$.  A planet at
\textit{tropical} $\lambda = 30\degree$ (the start of Taurus in
the tropical frame) is located at approximately \textit{sidereal}
$\lambda = 6\degree$ (still in Aries in the sidereal frame).  All
eleven bodies of the \zh{七政四餘} system shift by roughly one
zodiac sign under the change of convention.

This $\approx 24\degree$ discrepancy is the coordinate source of
the inter-system disagreements documented in the companion paper
\cite{liu2020missionary}.  When modern practitioners inadvertently
mix the Ming (sidereal) and Qing (tropical) conventions, a
systematic error of this magnitude is introduced into every
positional judgment.

\subsection*{Identity Thread: Ketu Restored, \zh{紫氣} Abolished}

The Jesuit reform was no less consequential on the identity thread.
Two changes stand out.

First, \zh{計都} reverted to the descending lunar node.  The
apogee reidentification that had persisted from the Tang period
through the Song and Ming codifications---a reidentification
demonstrated by Niu Weixing's computational analysis
\cite{niu1995ketu}---was reversed.  In the Schall reform,
\zh{計都} is once again the descending node, its original Indian
identity restored after approximately 900 years in the guise of
the apogee.  The reversal was not the result of Chinese scholarly
debate about the original Indian sources; it was a consequence of
the imported European computational framework, which had no concept
of an ``apogee body'' and naturally assigned the node bodies their
standard Indian identities.

Second, \zh{紫氣} was abolished entirely.  The purely Chinese
fictitious point with its 28-year cycle had no place in the
European-derived framework.  Its removal was a decisive reduction of
the tradition's indigenous Chinese additions, eliminating the one
body that had been added by Chinese astronomers with no foreign
precedent.

An elemental inconsistency followed from the \zh{計都} reidentification.
Within the \zh{五行} (\textit{w\v{u}x\'{\i}ng}, ``five phases'')
assignments that govern astrological interpretation, the Cai Boli
(\zh{蔡伯勵}) tradition---which preserves the Qing convention and
remains influential in Cantonese-speaking communities---encodes a
fire/earth swap relative to the earlier convention.  Because
\zh{計都} as apogee and \zh{計都} as descending node occupy
different positions in the sky at the same moment, their
phase-element affinities, derived from their positions in the
annual cycle, differ systematically.  The resulting
\zh{計北羅南}/\zh{計南羅北} toggle (``Ketu north, Rahu south'' vs.\
``Ketu south, Rahu north'') that appears in modern software
encoding reflects not a computational error but a genuine dual
inheritance from the pre- and post-Jesuit conventions
\cite{liu2020missionary}.

\subsection*{The Calendar Case (1664--1669)}

The Jesuit reform did not go unchallenged.  In 1664, Yang Guangxian
(\zh{楊光先}), a Han Chinese official and amateur astronomer,
submitted a memorial to the Kangxi Emperor accusing the Jesuits
of errors in their astronomical methods, of promoting heterodox
cosmology (the heliocentric model was philosophically suspect in
the Confucian framework), and of using their calendar monopoly to
subvert Chinese institutions \cite{jami2015calendar}.

The accusation had political as well as astronomical dimensions.
The Kangxi Emperor was a child; the regency was dominated by
conservative bannermen suspicious of European influence.  In 1665
Schall was convicted, sentenced to death, and subjected to the
\textit{lingchi} (``death by a thousand cuts'').  The sentence was
commuted---according to contemporary accounts, because a severe
earthquake struck Beijing shortly after sentencing, interpreted as
a sign of heaven's displeasure at the verdict.  Schall died under
house arrest later that year.  Yang Guangxian was appointed to lead
the Bureau of Astronomy, reverting briefly to Chinese computational
methods.

The scientific resolution came in 1669.  The Kangxi Emperor, now
governing in his own right and genuinely curious about mathematics
and astronomy, arranged a public astronomical competition.  Ferdinand
Verbiest (\zh{南懷仁}), a Flemish Jesuit who had served under
Schall, challenged Yang Guangxian to predict the length of a shadow
cast by a gnomon at noon on a specified date, the position of the
Sun, Moon, and visible planets for a specified day, and the timing
of an impending lunar eclipse.  Verbiest's predictions, based on
European methods, were consistently more accurate.  Yang Guangxian
conceded, resigned his post, and was escorted home in disgrace
\cite{jami2015calendar}.

Verbiest was appointed to lead the Bureau of Astronomy.  The
European astronomical infrastructure was reinstated.  The Qing
\zh{七政四餘} convention---ecliptic tropical coordinates, Ketu as
descending node, no \zh{紫氣}---was confirmed as the official
standard.

\subsection*{The Kangxi Synthesis and Pragmatic Dualism}

The Kangxi Emperor's resolution of the Calendar Case was not merely
a vindication of European methods.  It was the articulation of a
pragmatic dualism that persists in the tradition to this day: Western
astronomical data for computational precision, Chinese astrological
framework for interpretive content.

The Emperor did not mandate that Chinese practitioners adopt European
cosmology.  He mandated that official calendrical computation use
European positional astronomy.  The astrological interpretations
attached to planetary positions---the \zh{五行} element assignments,
the house-by-house delineations, the fate-technical categories
inherited from the medieval Chinese synthesis---remained Chinese.
The result was a hybrid of a different kind from the Ming
\zh{似黃道恆星製}: not a mathematical hybrid of coordinate systems
but an institutional hybrid of computational authority and
interpretive sovereignty \cite{liu2020missionary,jami2015calendar}.

\subsection*{Significance: The Ironic Circle}

The Jesuit disruption completes what may be called the coordinate
irony of \zh{七政四餘} history.  The tradition began with ecliptic
coordinates imported from India.  Over six centuries it migrated
away from the ecliptic, first to the equator (Yuan), then to the
quasi-ecliptic hybrid (Ming), achieving maximum sinicization.  The
Jesuit reform pulled it back to the ecliptic plane---but not to the
Indian ecliptic.  The new ecliptic system had a European zero point
(vernal equinox) rather than a sidereal Indian zero point, making
it tropical rather than sidereal.  The coordinate trajectory
returned to its plane of origin but not to its origin.

The contrast with Mode~3 (indigenous restructuring) illuminates what
makes Mode~5 distinctive.  The Yuan reform was driven by Chinese
astronomers applying the tools of Chinese observational astronomy to
an inherited system; the result was organic, internally motivated,
and absorbed over several generations.  The Qing reform was driven
by Jesuit astronomers with institutional backing from the imperial
court; the result was abrupt, externally imposed, and generated a
fractured dual tradition that survives into the present.

The fracture is visible in modern practice.  The Cai Boli tradition
(predominantly Cantonese, post-Qing) uses tropical coordinates and
identifies \zh{計都} as the descending node.  The mainstream
\zh{果老星宗} tradition (drawing on Ming texts) uses sidereal
coordinates and identifies \zh{計都} as the apogee.  Both claim
historical legitimacy.  Both are historically legitimate.  The
$\approx 24\degree$ positional discrepancy between them is not an
error; it is the fossil record of 1645.

% === §8 Algorithmic Fossilization =============================================
\section{Algorithmic Fossilization}\label{sec:fossil}
% ~1,200 words — Task 10

\textit{Section placeholder.}

% === §9 Computational Verification ============================================
\section{Computational Verification}\label{sec:computation}
% ~1,500 words — Task 10

\textit{Section placeholder.}

% === §10 Discussion ===========================================================
\section{Discussion}\label{sec:discussion}
% ~1,200 words — Task 11

\textit{Section placeholder.}

% === §11 Conclusion ===========================================================
\section{Conclusion}\label{sec:conclusion}
% ~500 words — Task 11

\textit{Section placeholder.}

% === Bibliography =============================================================
\bibliographystyle{unsrtnat}
\bibliography{silk-road-astronomy}

\end{document}
